% META: {"id": "TSPE-LIMFCT-EX-049", "chapitre": "TSPE-LIMITES-FONCTIONS", "type_objet": "exercice", "capacites": ["TSPE-LIMFCT-C3"], "capacites_codes": ["C3"], "methodes": ["M3"], "parcours": 3, "competences": ["raisonner", "calculer", "communiquer"], "duree_min": 25, "mode_creation": "ex_nihilo", "sources_inspiration": [], "fichier_tex": "chapitres/TSPE-LIMITES-FONCTIONS/exercices/TSPE-LIMFCT-EX-049.tex", "corrige_tex": "chapitres/TSPE-LIMITES-FONCTIONS/corriges/TSPE-LIMFCT-CO-049.tex", "status": "generated"}
% BEGIN-VERIFY
% from sympy import *
% x = symbols('x')
% f = (x**2 - 2*x + 3)*exp(-x)
% assert limit(f, x, oo) == 0
% assert limit(f, x, -oo) == oo
% fp = diff(f, x)
% assert simplify(fp - (-x**2 + 4*x - 5)*exp(-x)) == 0
% END-VERIFY
\begin{exercice}{TSPE-LIMFCT-EX-049}{3}{25}
\textbf{(Type bac)} Soit $f(x) = (x^2 - 2x + 3)\,\mathrm{e}^{-x}$ definie sur $\mathbb{R}$.
\begin{enumerate}
  \item Determiner les limites de $f$ en $+\infty$ et $-\infty$. Interpreter graphiquement.
  \item Calculer $f'(x)$ et montrer que $f'(x) = (-x^2 + 4x - 5)\,\mathrm{e}^{-x}$.
  \item Etudier le signe de $f'(x)$.
  \item Dresser le tableau de variations complet de $f$.
\end{enumerate}
\end{exercice}
